% ============================================================
% ch02.tex —— 第 2 章 铺地砖问题与辗转相除
% ============================================================
\chapter{铺地砖问题与辗转相除}

上一章我们学会了"压扁"大数。这一章解决一个更古老的问题：两个数，怎么找到它们"最大的共同尺寸"？别被这句话绕晕，直接上工地。

\begin{kaiwei}
一间 $12$ 米长、$18$ 米宽的矩形展厅要铺正方形地砖，要求：地砖不许切割，恰好铺满，砖缝横平竖直。地砖的边长最大是多少？
\end{kaiwei}

\section{从试数到猜出规律}

先像工人一样试数，把每种边长都过一遍：

\textbf{边长 $1$ 米}：$12 \div 1 = 12$ 列，$18 \div 1 = 18$ 行，铺 $216$ 块。行是行，但砖太小，缝比砖还多，甲方不满意。

\textbf{边长 $2$ 米}：$12 \div 2 = 6$ 列，$18 \div 2 = 9$ 行，铺 $54$ 块，横竖都除尽。可以。

\textbf{边长 $3$ 米}：$12 \div 3 = 4$，$18 \div 3 = 6$，铺 $24$ 块。更好。

\textbf{边长 $4$ 米}：$12 \div 4 = 3$ 列没问题；但 $18 \div 4 = 4.5$——\textbf{横向铺不满}，最后剩半块的地方。不合格。

\textbf{边长 $5$ 米}：$12 \div 5$ 除不尽。不合格。

\textbf{边长 $6$ 米}：$12 \div 6 = 2$，$18 \div 6 = 3$，铺 $2\times 3 = 6$ 块，严丝合缝。漂亮。

\textbf{边长 $7$ 米及以上}：$7, 8, 9, \dots$ 都不是 $12$ 的因数（最多试到 $12$），全部出局。

答案是 $\boxed{6}$ 米。翻译成数学：我们要找的是"既能整除 $12$、又能整除 $18$ 的数"里最大的一个——也就是 $12$ 和 $18$ 的\Keyword{最大公因数}，记作 $\gcd(12, 18) = 6$。（$\gcd$ 是 greatest common divisor 的缩写。）

小数字这样试没问题，可要是展厅换成 $2026$ 米 $\times$ $1998$ 米呢？把两千多个因数一个个列出来？工程浩大，还容易漏。两千三百年前，欧几里得在《几何原本》第七卷里给出了一个快得离谱的机器，两千年后的计算机教材第一章几乎都拿它开刀。先看它运转，再问为什么。

\section{欧几里得的机器}

\begin{dingli}[欧几里得算法（辗转相除法）]
求 $\gcd(a, b)$（$a \ge b > 0$）：用 $b$ 除 $a$，记余数为 $r_1$；再用 $r_1$ 除 $b$，记余数为 $r_2$；再用 $r_2$ 除 $r_1$……如此"辗转"，余数一路变小，直到某次余数为 $0$。\textbf{最后一个非零的除数就是 $\gcd$}。
\end{dingli}

看它碾过 $\gcd(2026, 1998)$：
\[
\begin{array}{rcl}
2026 & = & 1 \times 1998 + \mathbf{28} \\
1998 & = & 71 \times 28 + \mathbf{10} \\
28 & = & 2 \times 10 + \mathbf{8} \\
10 & = & 1 \times 8 + \mathbf{2} \\
8 & = & 4 \times 2 + \mathbf{0} \quad\Longleftarrow\ \text{余数归零，刹车！}
\end{array}
\]
最后的非零除数是 $2$，所以 $\gcd(2026, 1998) = \boxed{2}$。四千年前的展厅，五次带余除法定案。规则就一句话：\textbf{旧除数当新的被除数，旧余数当新的除数，直到除尽为止}。

顺手在旁边记一笔（第 3 章要回来对账）：分解 $2026 = 2 \times 1013$、$1998 = 2 \times 3^3 \times 37$，两个"素数配方"的公共部分恰好是 $2$——两种方法殊途同归。

\section{为什么裁到最后剩下的就是答案}

机器快是快，凭什么？用铺地砖的眼光看一遍，这是最直观的证法（建议真的画一张 $12\times 18$ 的图，边看边比划）：

\textbf{第一步：把大问题裁小。}想用边长 $d$ 的方砖铺满 $12\times 18$ 的地面。先在左端铺一个 $12\times 12$ 的正方形区域（用 $d$ 能整除 $12$，所以这个正方形铺得严丝合缝），右边剩下一块 $12 \times (18-12) = 12\times 6$ 的长条。\textbf{关键来了：}$d$ 能铺满 $12\times18$，当且仅当 $d$ 能铺满剩下的 $12\times 6$ 长条。（"当且仅当"是数学的严丝合缝说法：两面都成立。）于是"$d$ 是 $18$ 和 $12$ 的公因数"与"$d$ 是 $12$ 和 $6$ 的公因数"是\textbf{同一件事}——公因数一个没多、一个没少，难题被裁小了。

\textbf{第二步：反复裁。}$12 \times 6$ 的长条里，裁掉两个 $6\times 6$（$12 = 2\times 6$），剩下 $0$。也就是说问题一路缩小成"$d$ 是 $6$ 和 $0$ 的公因数"——任何数都是 $0$ 的因数（$0 = d\times 0$），所以这时"$d$ 是 $6$ 的因数"。最大的 $d$ 就是 $6$。答案从裁剪过程里自己长出来了。

对照带余除法的链条：
\[
18 = 1\times 12 + 6,\qquad 12 = 2\times 6 + 0 .
\]
每一步"旧长条裁掉正方块、剩新长条"，正是"旧数除以旧余数"。裁剪必停（长条越来越短，最短到零），停下的地方就是答案。\textbf{辗转相除的每一步，就是"用小方砖先铺掉能铺的部分，把问题裁小一圈"}；余数就是裁剩下的边角。

用代数把同一段话写死（严格版）：若 $a = bq + r$，则"$d$ 整除 $a$ 且 $d$ 整除 $b$"$\iff$"$d$ 整除 $b$ 且 $d$ 整除 $r$"。
- 往右：$d\mid a$，$d\mid b$，则 $r = a - bq$ 是两个 $d$ 的倍数之差，故 $d \mid r$；
- 往左：$d\mid b$，$d\mid r$，则 $a = bq + r$ 是两个 $d$ 的倍数之和，故 $d\mid a$。

所以 $(a, b)$ 与 $(b, r)$ 拥有\textbf{完全相同的公因数集合}，最大公因数自然相等：$\gcd(a,b) = \gcd(b,r)$。链条一路替换下去，余数一路变小，最后一步 $\gcd(d, 0) = d$。

\begin{faming}
\textbf{土名}：辗转相除、互相裁边角 $\to$ \textbf{正式名}：欧几里得算法（Euclidean algorithm）。它是人类有文字记载以来\textbf{第一个成型的算法}——不光给出答案，还给出"照着做必得答案"的机械流程。两千年来它稳坐"最古老且仍在前线服役的算法"宝座：今天的计算机求最大公因数、加密软件算密钥，跑的还是这几行。
\end{faming}

\section{意外收获：最大公因数是"调"出来的}

辗转相除还藏着一张隐藏底牌，第 6 章造密码锁、第 10 章发"除法通行证"都要用到它。回头看 $\gcd(387, 210)$ 的链条：
\[
\begin{array}{rcl}
387 & = & 1\times 210 + 177 \\
210 & = & 1\times 177 + 33 \\
177 & = & 5\times 33 + 12 \\
33 & = & 2\times 12 + 9 \\
12 & = & 1\times 9 + 3 \\
9 & = & 3\times 3 + 0
\end{array}
\]
答案是 $3$。现在做一件奇怪的事：\textbf{从最底下往回倒着代}。最后第二行说 $3 = 12 - 1\times 9$；而 $9$ 又来自上一行，$9 = 33 - 2\times 12$，代入：
\[
3 = 12 - 1\times(33 - 2\times 12) = 3\times 12 - 1\times 33 .
\]
再往上，$12 = 177 - 5\times 33$，代入：
\[
3 = 3\times(177 - 5\times 33) - 33 = 3\times 177 - 16\times 33 .
\]
再往上，$33 = 210 - 177$：
\[
3 = 3\times 177 - 16\times(210 - 177) = 19\times 177 - 16\times 210 .
\]
最后，$177 = 387 - 210$：
\[
\boxed{3 = 19\times 387 - 35\times 210 .}
\]
验算：$19 \times 387 = 7353$；$35\times 210 = 7350$；差恰好 $3$。\checkmark

奇迹发生了：\textbf{最大公因数可以表示成原来两个数的整数组合}——不是碰巧，是因为链条上每一个余数都由前两个数"调"出来（$177 = 387 - 210$，$33 = 210 - 177$……每个数都是上面两个数的组合，组合的组合还是组合，一路倒回去就调出了 $\gcd$）。这条结论叫\Keyword{裴蜀等式}（Bézout，法国数学家，代表作是那本著名的《数学辞典》序言）：

\begin{dingli}[裴蜀等式]
对任何正整数 $a, b$，存在整数 $x, y$（可正可负）使
\[
ax + by = \gcd(a, b).
\]
特别地，若 $\gcd(a,b) = 1$（\Keyword{互素}），则存在整数 $x,y$ 使 $ax + by = 1$。
\end{dingli}

互素的关系有多常见？随手抓两把：$8$ 和 $9$ 互素（$8\times(-1) + 9\times 1 = 1$，一眼可验）；$2\times 3 - 1\times 5 = 1$，所以 $3$ 和 $5$ 互素；$2027$ 和 $2028$ 这种\textbf{相邻整数}永远互素——因为任何公因数都得整除它们的差 $1$，而整除 $1$ 的只有 $\pm 1$。相邻必互素，这是后面章节反复免费使用的小零件。

\section{互素的两个速判法}

判断互素不必每次都跑完整台机器，两条捷径：

\textbf{捷径一（减法版）}：$\gcd(a, b) = \gcd(a - b, b)$。因为同时整除 $a$ 和 $b$ 的数，必然整除差 $a - b$；反过来，整除 $a-b$ 和 $b$ 的数，也整除 $(a-b)+b = a$。两个问题的"公因数集合"一模一样。辗转相除本质上是这条捷径的高效连用（先用减法会太慢，除法是一口气减去很多次）。

\textbf{捷径二（素数版）}：$a$ 和 $b$ 互素 $\iff$ 它们\textbf{没有共同的素因数}。比如 $2026 = 2\times 1013$、$1998 = 2\times 3^3\times 37$，公共素因数只有 $2$，所以 $\gcd = 2$。而 $\gcd$ 的"素数配方"恰好是两者配方中每个素数取较小次幂：$\gcd(2^1\cdot 1013^1,\ 2^1\cdot 3^3\cdot 37^1) = 2^1 = 2$。

捷径二已经嗅到下一章的气息：\textbf{素数是整数的原子，一切因数问题最终都是原子问题}。

\begin{lishi}
欧几里得生活的年代（约公元前 330--前 275）比秦始皇统一六国还早一百年。他编写的《几何原本》十三卷，把当时希腊数学的成果按"定义—公设—定理"的骨架重新组织，成为历史上再版次数第二多的书（第一是《圣经》）。第七到第九卷讲的是数论，我们这两章的内容——辗转相除、素数、唯一分解、素数无穷——全部来自那里。也就是说：你今天在草稿纸上做的辗转相除，和一个两千三百年前的希腊人用的是同一台机器，连步骤都一样。
\end{lishi}

\begin{dongshou}
\begin{itemize}
  \yisi 用辗转相除求 $\gcd(120, 45)$，并画出每一步"裁边角"的矩形示意图（$120 = 2\times45 + 30$，然后 $45 = 1\times 30 + 15$，然后 $30 = 2\times 15 + 0$）。
  \yier 求 $\gcd(2026, 1998)$ 的裴蜀表示：找整数 $x, y$ 使 $2026x + 1998y = 2$。（照着 2.4 节"倒着代入"的流程一步步走。给你对账用的中间结果：倒代到最后一步会得到 $2 = 3\times 1998 - 214\times 2026 + 214\times 1998$，合并同类项即得答案；验算 $217\times 1998 - 214\times 2026 = 433566 - 433564 = 2$。\checkmark 这题的价值全在"倒代"的手感，别跳步。）
  \yier 证明捷径一：若 $d$ 整除 $a$ 和 $b$，则 $d$ 整除 $a - b$；再说明反过来为什么也对（若 $d$ 整除 $a - b$ 和 $b$，则 $d$ 整除 $(a-b)+b = a$）。
  \yier 利用"相邻整数互素"秒答：$\gcd(n, n+1)$ 是多少？$\gcd(n!, (n!+1))$ 呢？（第二个答案为 $1$——它将出现在第 4 章构造"合数荒漠"的证明里，先在这里埋好。）
\end{itemize}
\end{dongshou}

\begin{shaonao}
连续 $k$ 个整数 $n, n+1, \dots, n+k-1$ 的公因数 $d$，必然整除其中任何两个的差——而任何两数之差都不超过 $k-1$，更妙的是 $d$ 整除 $(n+1) - n = 1$ 吗？不一定（除非 $k \ge 2$ 且取相邻对）。请把这条一般规律证出来：\textbf{$\gcd(n, n+k)$ 整除 $k$}。（把 $n + k$ 减去 $n$，用捷径一。）然后用它三秒判断：$\gcd(1000, 1005)$ 只可能是 $5$ 的因数；$\gcd(2^{100}, 2^{100}+6)$ 整除 $6$，又两数都是偶数，所以 $\gcd$ 恰是 $2$ 还是 $6$？算一下 $2^{100}$ 除以 $3$ 的余数来定案（用第 1 章的周期表思路：$2 \equiv -1 \pmod 3$，所以 $2^{100} \equiv 1 \pmod 3$，于是 $2^{100} + 6 \equiv 1 + 0 = 1 \pmod 3$，$3$ 不整除它——$\gcd = 2$）。
\end{shaonao}

\begin{chengshi}
裴蜀等式的完整证明（对辗转相除的步数做归纳，"每个余数都是 $a,b$ 的组合"从链条底部归纳上去）我们没有形式化展开，只演示了"倒着代入"的操作——操作本身就是证明的灵魂：归纳的每一步就是"再代入一次"。想要教科书版本，任何一本初等数论教材第一节都有，两行。另外欧几里得算法"很快"有精确版本（拉梅定理：步数不超过较小数的十进制位数的五倍），超纲，略。
\end{chengshi}

下一章，我们请出整数世界的"元素周期表"：素数。你会看到本章的互素与裴蜀，在"唯一分解"的大定理里当关键零件。
